\documentclass{article}
\usepackage{graphicx}
\usepackage{amsmath}
\usepackage{amssymb}
\usepackage{pgffor}

\foreach \x in {A,B,C,D,E,F,G,H,I,J,K,L,M,N,O,P,Q,R,S,T,U,V,W,X,Y,Z}{
  \expandafter\xdef\csname mc\x\endcsname{\noexpand\mathcal{\x}}
}

\newcommand{\N}{\mathbb{N}}
\newcommand{\Z}{\mathbb{Z}}
\newcommand{\Q}{\mathbb{Q}}
\newcommand{\R}{\mathbb{R}}
\newcommand{\C}{\mathbb{C}}
\newcommand{\E}{\mathbb{E}}
\renewcommand{\P}{\mathbb{P}}

\newcommand{\seq}{\mathrm{SEQ}}
\newcommand{\set}{\mathrm{SET}}
\newcommand{\cyc}{\mathrm{CYC}}

\newcommand{\bigp}[1]{\left(#1\right)}
\newcommand{\eps}{\varepsilon}

\newcommand{\problem}[1]{\section*{Problem #1}}
\renewcommand{\part}[1]{\subsection*{Part (#1)}}


\title{Generating Functions 6}
\author{Siddharth K}
\date{October 2026}

\begin{document}

\maketitle

\problem{2}
We'll show that there is some driver that cannot park iff there exists $i$ such that $|f^{-1}([i])| < i$ (there exists a parking spot $i$ that fewer than $i$ cars prefer). First, assume the latter. Then $> n - i$ have to park in the spots beyond $i$. There are $n - i$ of them, so you cannot park $> n - i$ people in them, so at least one cannot park.

For the other direction, assume that the car that wanted the $i$th spot cannot park, which means that all spots from $i$ to $n$ are taken. Say $j$ is the highest available spot, so that $j < i$. Then $> n - j + 1$ people wanted some spot between $j + 1$ and $n$, so the number of people who want a spot between $1$ and $j$ is $\le n - (n - j + 1) = j - 1 < j$, as desired.

\problem{3}
Suppose we have $n$ cars who want to park in $n + 1$ spaces, but the cars can wrap around back to spot $1$ if they don't find a spot between their preference and $n + 1$. Then every car can park, and there will be one empty spot. Consider the arrangements where that empty spot is $n + 1$: every car's preferred spot is in $[n]$, and no cars wrap around past $n + 1$ (because then it would be occupied). Therefore all cars park in a way that satisfies the original parking problem. The number of such arrangements is
\[
    \frac{(n + 1)^n}{n + 1} = (n + 1)^{n - 1}.
\]

\problem{4}
Given a parking function $f$, we can form a tree on $\{0, 1, \ldots, n\}$ in this way: for each $i \in n$,

\begin{enumerate}
    \item if $f(i) = 1$, connect $i$ to $0$,
    \item otherwise, connect $i$ to the car occupying the spot before car $i$'s desired spot: $f(i) - 1$.
\end{enumerate}

This is a tree because for each $i$, its parent is either $0$ or occupies the spot $f(i) - 1$, which is always less than the spot that car $i$ ended up at. So following each parent follows a path of strictly decreasing parking spots until vertex $0$. This means there are no cycles, so it is a tree.

We can reverse this process by defining $f$ as
\[
    f(i) = (\text{parking spot that the parent of } i \text{ ended at}) + 1.
\]

\problem{5}
Given a rooted forest on the set $[n]$, we can construct one tree by adding a vertex labeled $0$ and connecting it to the root of each tree in the forest. The resulting tree has $n + 1$ vertices. To reverse this process, take a tree on the set $\{0, 1, \ldots, n\}$. Then remove $0$, take the resulting components as the forest, and assign the root of each tree to be the vertex that was connected to $0$ in the original tree. By Cayley's theorem, the number of trees on $\{0, 1, \ldots, n\}$ is $(n + 1)^{n - 1}$, so that is also the number of rooted forests on $[n]$.

\problem{6}
Given a parking function on $[n]$, make an unrooted labeled tree on $n + 1$ vertices using the bijection from problem 4, and then make a rooted forest from that tree using the bijection from problem 5 (and similarly in reverse).

\problem{8}
We can define a labeled hierarchy recursively. Given a labeled hierarchy, there are two possibilities: either the tree consists of one labeled leaf, or the tree has an unlabeled root with at least $2$ children. So the specification is
\[
    \mcL \cong \mcZ + \set_{\ge 2}(\mcL).
\]
So the generating function $L(z)$ satisfies
\[
    L(z) = z + e^{L(z)} - 1 - L(z),
\]
or
\[
    L(z) = \frac{1}{2}(z + e^{L(z)} - 1).
\]
I don't think it's possible to get an explicit formula for $L(z)$ given this.

\problem{9}
This is the same as the last problem but the set must have size $2$, so the specification is
\[
    \mcM \cong \mcZ + \set_2(\mcM),
\]
and the generating function satisfies
\[
    M(z) = z + \frac{1}{2}M(z)^2,
\]
or
\[
    M(z)^2 - 2M(z) + 2z = 0.
\]
Isolating $M(z)$ gives
\[
    M(z) = 1 \pm \sqrt{1 - 2z}.
\]
Since $M(0) = 0$, we choose the negative solution:
\[
    M(z) = 1 - \sqrt{1 - 2z}.
\]
By the generalized binomial theorem,
\[
    (1 - 2z)^{1/2} = \sum_{n = 0}^\infty \binom{1/2}{n}(-2z)^n.
\]
So for $n \ge 1$, the answer is
\[
    m_n = -n!\binom{1/2}{n}(-2)^n.
\]

\problem{11}
Given a function $f: [n] \to [n] \cup \{\perp\}$, consider its functional graph. Since some vertices map to $\perp$, it consists of several connected components where each component is either:
\begin{enumerate}
    \item a rooted tree whose root maps to $\perp$, or
    \item a directed cycle of rooted trees.
\end{enumerate}
So the specification for $\mathcal{PF}$ is
\[
    \mathcal{PF} \cong \set(\mcT + \cyc(\mcT)),
\]
where $\mcT$ is the class of labeled rooted trees. Therefore, the generating function is
\[
    PF(z) = e^{T(z)} \cdot e^{\log \frac{1}{1 - T(z)}} = \frac{e^{T(z)}}{1 - T(z)}.
\]
Since
\[
    T(z) = ze^{T(z)},
\]
we can differentiate both sides to get

\begin{align*}
    T'(z) &= e^{T(z)} + ze^{T(z)}T'(z) \\
    &= e^{T(z)} + T(z)T'(z), \\
    T'(z)(1 - T(z)) &= e^{T(z)}, \\
    T'(z) &= \frac{e^{T(z)}}{1 - T(z)},
\end{align*}

which is the same as earlier, so
\[
    PF(z) = T'(z).
\]
By Cayley's theorem, $t_n = n^{n - 1}$, so

\begin{align*}
    T(z) &= \sum_{n \ge 1} n^{n - 1} \frac{z^n}{n!}, \\
    PF(z) = T'(z) &= \sum_{n \ge 1} n^{n - 1} \frac{z^{n - 1}}{(n - 1)!} \\
    &= \sum_{n \ge 0} (n + 1)^n \frac{z^n}{n!}.
\end{align*}
Comparing coefficients gives
\[
    pf_n = (n + 1)^n,
\]
which matches with the straightforward counting argument: $n + 1$ output options for each of the $n$ possible inputs.

\problem{13}
Given a permutation with $k$ minima, we can split the permutation before each minimum (except the first one), leaving $k$ remaining blocks. Then use those blocks as cycles for a new permutation. To go in reverse, take the $k$ cycles of the permutation, write the cycle so that the smallest element appears first, and then concatenate them in order of ascending first elements. The result is a permutation where each start of a cycle is now a minimum.

\end{document}